\section{Introduction}
\label{sec:intro}

Adding two steering vectors and injecting the sum into a language model's
residual stream is one of the simplest interventions in mechanistic
interpretability. Practitioners use this trick to combine behavioral
modifications---\eg making a model simultaneously more formal \emph{and} more
concise---under the implicit assumption that linear directions compose via
vector addition. But do they?

The linear representation hypothesis provides theoretical grounding: concepts
are encoded as directions in activation space, and causally separable concepts
occupy orthogonal subspaces~\citep{park2023linear}. Under this view, adding two
orthogonal concept vectors should steer both concepts independently.
\citet{park2024geometry} extend this to show that hierarchically related
concepts compose via vector arithmetic, and \citet{todd2023function} demonstrate
that in-context-learned function vectors can be composed when sub-operations are
separable. Yet these results cover isolated cases. No study has systematically
tested \emph{which} linear directions compose and \emph{what predicts} their
composability across a diverse set of concept types.

This gap has practical consequences. Activation steering is increasingly used
for model control~\citep{stolfo2024improving}, but individual steering vectors
are already unreliable---per-sample steerability is highly
variable~\citep{tan2024analyzing}, and single directions activate unintended
features~\citep{chalnev2024improving}. Composing vectors without understanding
their geometric relationship risks compounding these failure modes.

We address this gap with a systematic empirical study. We extract linear
directions for 10 binary concepts spanning 5 semantic categories in \pythia
(32~layers, $\dmodel = 2560$) and test all 10 pre-selected concept pairs for
compositionality using three complementary metrics: (1)~pairwise cosine
similarity, measuring geometric alignment; (2)~probing-based preservation,
measuring whether both concepts remain linearly decodable from the composed
direction; and (3)~steering-based component preservation, measuring whether the
composed vector shifts model logits in both intended directions. Our main
findings are:

\begin{itemize}[leftmargin=*,itemsep=0pt,topsep=0pt]
    \item We show that most mean-difference concept directions are
    \textbf{surprisingly non-orthogonal} (median $\abscossim = 0.91$ at
    layer~20), even for seemingly unrelated concepts, challenging the
    assumption that independent concepts produce independent directions.
    \item We demonstrate that \textbf{probing-based composition metrics are
    misleadingly optimistic}: preservation is uniformly high (${\sim}0.95$),
    while steering reveals stark composition failures for aligned direction
    pairs.
    \item We identify \textbf{absolute cosine similarity} between direction
    pairs as the strongest predictor of steering composition success, with
    near-orthogonal pairs ($\abscossim < 0.3$) composing well and
    anti-parallel pairs failing catastrophically.
    \item We provide the first \textbf{systematic taxonomy} of
    compositionality across concept categories, finding that hierarchical
    pairs compose best and entangled pairs compose worst, consistent with
    theoretical predictions from \citet{park2023linear}.
\end{itemize}
